\subsection{EMBEDDING OF A LOCALLY COMPACT SPACE IN A COMPACT SPACE}

THEOREM 4 (Alexandroff). If X is any locally compact space, there exists a compact space \(X'\) and a homeomorphism f of X onto the complement of a point in \(X'\) . Furthermore, if \(X_{1}'\) is another compact space such that there is a homeomorphism \(f_{1}\) of X onto the complement of a point in \(X_{1}'\) , then there is a unique homeomorphism g of \(X'\) onto \(X_{1}'\) such that \(f_{1}=g\circ f\) .

Let us begin by proving the second assertion of the theorem. Let

\[
f (\mathbf {X}) = \mathbf {X} ^ {\prime} - \left\{\omega \right\} \quad \text { and } \quad f _ {1} (\mathbf {X}) = \mathbf {X} _ {1} ^ {\prime} - \left\{\omega_ {1} \right\}.
\]

If the homeomorphism \(g\) exists it must be unique, for by definition we have \(g(x') = f_1(\overline{f}^1 (x'))\) if \(x' \neq \omega\) and therefore \(g(\omega) = \omega_1\) . It remains to show that the bijection \(g: \mathbf{X}' \to \mathbf{X}_1'\) thus defined is bicontinuous; since \(\mathbf{X}'\) and \(\mathbf{X}_1'\) are interchangeable we need only show that the image under \(g\) of a neighbourhood of a point \(x' \in \mathbf{X}'\) is a neighbourhood of \(g(x')\) in \(\mathbf{X}_1'\) . This is obvious from the definition of \(g\) if \(x' \neq \omega\) . If \(x' = \omega\) , let \(V'\) be an open neighbourhood of \(\omega\) in \(X'\) ; then \(X' - V' = K\) is closed in \(X'\) and therefore compact (no. 3, Proposition 3) and is contained in \(f(X)\) ; hence \(g(K) = f_1(\overline{f}^1 (K))\) is compact (no. 4, Theorem 2, Corollary 1). It follows that \(g(V') = X_1' - g(K)\) is an open neighbourhood of \(\omega_1\) (no. 3, Proposition 4). Hence \(g\) is a homeomorphism.

To prove the first part of the theorem, let \(X'\) be a set which is the sum of X and a set consisting of a single point \(\omega\) . We define a topology on \(X'\) by taking the set D of open subsets of \(X'\) to consist of all open subsets of X and all subsets of the form \((\mathbf{X}-\mathbf{K})\cup\{\omega\}\) , where K is a compact subset of X. Since any intersection of compact subsets of X is compact (no. 3, Propositions 3 and 4) and since any closed subset of a compact set is compact (no. 3, Proposition 3), it follows that \(\mathfrak{D}\) satisfies axiom \((\mathrm{O_I})\) ; and since any finite union of compact subsets of X is compact (no. 3, Proposition 5), \(\mathfrak{D}\) also satisfies \((\mathrm{O_{II}})\) . Every compact subset of X is closed in X (no. 3, Proposition 4) and therefore the topology induced on X by that of \(\mathbf{X}'\) is the original topology on X. Thus it remains to show that \(\mathbf{X}'\) is compact. In the first place, \(\mathbf{X}'\) is Hausdorff. For if \(x, y\) are any two distinct points of X, they have disjoint open neighbourhoods V, W respectively in X, and V, W are open in \(\mathbf{X}'\) ; on the other hand each \(x \in \mathbf{X}\) has a compact neighbourhood K in X, which is also a neighbourhood of \(x\) in \(\mathbf{X}'\) , while

\[
(\mathbf {X} - \mathbf {K}) \cup \{\omega \}
\]

is a neighbourhood of \(\omega\) in \(\mathbf{X}'\) and manifestly does not meet \(\mathbf{K}\) . Finally, \(\mathbf{X}'\) is quasi-compact. Let \((\mathbf{U}_{\lambda})_{\lambda \in \mathbf{L}}\) be an open covering of \(\mathbf{X}'\) ; then for at least one index \(\mu \in \mathbf{L}\) we have \(\mathbf{U}_{\mu} = (\mathbf{X} - \mathbf{K}_{\mu}) \cup \{\omega\}\) where \(\mathbf{K}_{\mu}\) is a compact subset of \(\mathbf{X}\) . Hence there is a finite subset \(\mathbf{H}\) of \(\mathbf{L}\) such that the sets \(\mathbf{U}_{\lambda}\) (for \(\lambda \in \mathbf{H}\) ) cover \(\mathbf{K}_{\mu}\) ; if \(\mathbf{J} = \mathbf{H} \cup \{\mu\}\) , then \((\mathbf{U}_{\lambda})_{\lambda \in \mathbf{J}}\) is a finite open covering of \(\mathbf{X}'\) , and the proof is complete.

Notice that if X is already compact, then \(\omega\) is an isolated point of the compact space \(X'\) ; hence \(X'\) is the sum ( \(\S\) 2, no. 4, Example 3) of the space X and the space \(\{\omega\}\) .

When a compact space \(X'\) has been constructed as above from a locally compact space X by adjoining an element \(\omega\) , it is often said that \(\omega\) is the ``point at infinity'' of \(X'\) , and that \(X'\) is obtained from X by adjoining a point at infinity. \(X'\) is also called the Alexandroff compactification or the one-point compactification of the locally compact space X.

* Example. If we apply Alexandroff's theorem to the real plane \(\mathbf{R}^2\) , we get a compact space homeomorphic to the sphere \(\mathbf{S}_2\) whose equation is \(x_1^2 + x_2^2 + x_3^2 = 1\) in \(\mathbf{R}^3\) . A homeomorphism of these two spaces may be described as follows: the point \(\omega\) (the point at infinity) adjoined to \(\mathbf{R}^2\) is mapped to \((0, 0, 1) \in \mathbf{S}_2\) , and every point \((x_1, x_2)\) of \(\mathbf{R}^2\) is mapped to the point where the line joining the points \((0, 1, 1)\) and \((x_1, x_2, 0)\) in \(\mathbf{R}^3\) meets \(\mathbf{S}_2\) again. This homeomorphism is known as stereographic projection.

